\chapter{Introduction} \label{chap:intro}

\section{Background and Significance}
Machine learning models deployed in production environments operate under the assumption that future data will resemble training data. This assumption is routinely violated: the distribution of the input features, the label distribution, or the conditional relationship between inputs and outputs changes over time, a phenomenon usually termed \textit{concept drift} or dataset shift \citep{Lu2018, Gama2014, Webb2016}. Left undetected, drift degrades model performance silently, accumulating technical debt and operational risk \citep{Sculley2015, Bayram2022}. \citet{WidmerKubat1996} argued that effective learning in non-stationary environments requires detecting context changes and recovering quickly from them.

The practical consequences are well documented. \citet{Vela2022} showed that standard machine learning models in four industries, including hospital operations, had growing error as the time since training increased in 91\% of 128 model-dataset pairs, even on data chosen to contain no abrupt changes. \citet{NestorEtAl2019} found that models for common clinical tasks failed to generalise across temporally shifted patient cohorts within the same hospital system. In such settings, an alarm that only states that drift occurred leaves practitioners to work out the cause, which in production systems can lie in upstream data producers, feature pipelines or changes in the external world \citep{Sculley2015}.

This thesis studies \textit{feature drift}, a change in the input distribution $P(X)$, which can be observed without labels. Labels in data streams often arrive late or not at all, so supervised detectors that monitor the error rate of a model, such as the Drift Detection Method (DDM) and the Early Drift Detection Method (EDDM) \citep{Gama2004, BaenaGarcia2006}, cannot always be used \citep{Lu2018, Lukats2024}. Unsupervised detectors instead compare a reference sample of the stream with recent data. Statistical methods apply two-sample tests such as the Kolmogorov--Smirnov (KS) test or the Maximum Mean Discrepancy (MMD) \citep{Rabanser2019, Gretton2012}, and adaptive windowing methods such as Adaptive Windowing (ADWIN) adjust the length of the reference window \citep{BifetGavalda2007}. Univariate tests are run once per feature and raise more false alarms as the number of features grows unless the significance level is corrected \citep{Rabanser2019}.

The discriminative approach recasts distribution comparison as a classification problem. \citet{LopezPazOquab2017} formalised it as the \textit{classifier two-sample test} (C2ST): if a classifier can separate reference data from current data based on the features alone, the two distributions differ. In machine learning practice the same procedure is called \textit{adversarial validation} \citep{Pan2020}. This thesis studies adversarial validation as a technique for feature drift detection and uses the term C2ST when it refers to the statistical test. \citet{Gozuacik2019} adapted this principle to data streams as the Discriminative Drift Detector (D3), which is one of the three detectors recommended by the benchmark of \citet{Lukats2024}. D3 trains a logistic regression classifier to separate the most recent window from the window just before it and signals drift when the area under the receiver operating characteristic curve (AUC) of the classifier reaches 0.7.

\section{Problem Statement}
Three properties of sliding-reference C2ST detectors such as D3 limit their use for unsupervised feature drift detection.

First, the reference window slides with the stream. When a drift is slow, as in gradual or incremental drift, the reference and the current window change together, the AUC stays low, and the drift can pass without an alarm. A reference that keeps older data would see the accumulated change, but it must not include the samples it is compared with.

Second, a fixed AUC threshold is not a calibrated statistical test. Permutation tests calibrate such a statistic without distributional assumptions \citep{Good2005}, but they assume that the samples are exchangeable under the null hypothesis. Real streams are often autocorrelated: consecutive samples are similar even when the distribution does not change, so a recent window differs from a reference by chance more often than shuffled data suggest. On such streams both a fixed threshold and a shuffling test raise false alarms although nothing has drifted. Permutation tests are also costly, because each permutation retrains the classifier; sequential Monte Carlo methods reduce the number of permutations a test needs \citep{Gandy2009}.

Third, D3 discards the classifier after computing the AUC, so an alarm does not say which features changed. SHapley Additive exPlanations (SHAP) \citep{Lundberg2017} and its exact tree variant TreeSHAP \citep{LundbergInteraction2020} have been applied to domain classifiers on batch data, where the ranking pointed to the shifted features \citep{Maggio2021}, and to the explanations of deployed models to detect shifts that change model behaviour \citep{Mougan2025}. Other feature-level methods need ground-truth labels and a trained model \citep{TRIPODD2024, HDPD2024}, or define feature-wise drift formally without a streaming detector \citep{HinderFeature2024}. Whether the SHAP attribution of a classifier-based detector on a data stream ranks the truly drifting features first, at the moment it raises an alarm, has to our knowledge not been measured against known ground truth.

\section{Research Objectives}
The thesis proposal set the goal of investigating adversarial validation as a feature drift detection technique. The objectives of this thesis are:

\begin{enumerate}
    \item To develop the Explainable Adversarial Drift Detector (EADD), an unsupervised C2ST for data streams that combines a reservoir-sampled reference window \citep{Vitter1985}, a permutation test whose null distribution comes from the stream's own history, an adaptive sequential permutation test (ASPT) and a TreeSHAP feature ranking after each alarm.

    \item To compare EADD with D3 and with D3 using the same Light Gradient Boosting Machine (LightGBM) classifier \citep{Ke2017} (D3-LightGBM) in terms of recall, mean time to detection and false alarms, on streams with known drift points and on streams without drift, including autocorrelated streams.

    \item To determine whether ASPT gives the same detections as a full permutation test with fewer permutations.

    \item To determine whether the TreeSHAP attribution of a classifier-based drift detector identifies the features that actually drifted.
\end{enumerate}

\section{Research Hypotheses}
\label{sec:intro-hypotheses}
The four hypotheses below are stated in their final form; Chapter~\ref{chap:methodology} gives the tests and decision rules, which were registered before the final runs. H1 was revised after the earlier experiments from all drift to slow drift, with the advisor's approval, and H4 was added.

% The hypothesis statements are defined once here and reused verbatim in Chapters 4 and 6.
\newcommand{\hypHone}{On streams with slow feature drift (gradual or incremental), EADD detects more drift points (H1a, higher recall) and detects them sooner (H1b, lower mean time to detection) than a sliding-reference C2ST with the same classifier (D3-LightGBM). On abrupt drift, its mean time to detection is at most 50 samples higher than that of D3-LightGBM (H1c, non-inferiority).}
\newcommand{\hypHtwo}{On streams without drift, where every alarm is false, EADD raises fewer false alarms than D3.}
\newcommand{\hypHthree}{ASPT gives the same recall, mean time to detection and false alarms as the full permutation test (H3a, equivalence) and uses fewer permutations (H3b).}
\newcommand{\hypHfour}{At the first detection of each drift point, the TreeSHAP feature attribution of a classifier-based drift detector ranks the truly drifting features above chance.}
\begin{description}
    \item[H1 (Detection of slow drift):] \hypHone
    \item[H2 (False alarms):] \hypHtwo
    \item[H3 (Sequential permutation test):] \hypHthree
    \item[H4 (Feature attribution):] \hypHfour
\end{description}
In H2, D3 is tested in two forms: D3-LightGBM, which uses EADD's classifier, is the primary comparison, and the original D3 with logistic regression is the secondary one. H4 concerns attribution in classifier-based detectors in general; it is measured with EADD's discriminator.

\section{Scope of the Study}
This thesis studies unsupervised detection of feature drift, that is, changes in $P(X)$, in tabular data streams with numerical features. Labels are never used for detection, so changes in $P(y \mid X)$ that leave $P(X)$ unchanged are outside the scope. The drift patterns studied are abrupt, gradual, incremental and recurring. The main evaluation uses 35 streams: 21 drift streams (seven each with abrupt, gradual and incremental drift, built from one synthetic generator and six real datasets into which drift was injected in known features), eight streams without drift and six autocorrelated streams without drift. Four real INSECTS streams with documented drift points \citep{Souza2020} and nine unannotated real streams from the benchmark of \citet{Lukats2024} are used in the wider benchmark. Feature attribution is scored only on synthetic and injected streams, where the drifting features are known. Drift adaptation, such as retraining the deployed model, is used only to score the unannotated streams and is not studied in itself.

\section{Assumptions}
The study relies on the following assumptions:
\begin{enumerate}
    \item Data seen since the last alarm come from the distribution that is considered normal, so that a difference between the reference and the current window indicates drift.
    \item Under the null hypothesis of no drift, earlier windows of the same stream behave like the current window, so that their AUCs against the reference form a valid null distribution. This replaces the exchangeability of individual samples, which autocorrelated streams violate.
    \item The drift points documented for the INSECTS streams \citep{Souza2020} and the injected drift points mark the true changes in the distributions.
\end{enumerate}

\section{Definitions of Terms}
\begin{description}
    \item[Feature drift:] a change over time in the distribution of the input features $P(X)$, also called covariate drift or virtual drift.
    \item[Concept drift:] in the narrow sense, a change over time in the conditional distribution $P(y \mid X)$; in the wider literature, any change in the joint distribution \citep{Gama2014, Lu2018}.
    \item[Classifier two-sample test (C2ST):] a test of whether two samples come from the same distribution, based on how well a classifier can tell them apart \citep{LopezPazOquab2017}.
    \item[Adversarial validation:] the practice of training a classifier to separate two data sets, such as training data and production data, and reading a high AUC as evidence that they differ \citep{Pan2020}. As a statistical test, it is a C2ST.
    \item[Reference window and current window:] the baseline sample $W_{ref}$ and the most recent sample $W_{cur}$ that a drift detector compares. A sliding reference holds the samples just before $W_{cur}$; a reservoir reference holds a uniform random sample of all data since the last alarm.
    \item[AUC:] the area under the receiver operating characteristic curve of the classifier that separates the two windows; 0.5 means that the windows cannot be separated.
    \item[Permutation test:] a test that compares the observed AUC with a null distribution of AUCs. In EADD the null AUCs come from earlier windows of the same stream (the history null).
    \item[Recall:] the share of drift points that are matched by a detection within a tolerance $T$ after them, with $T = 500$ samples for abrupt and recurring drift and $T = 1{,}500$ for gradual and incremental drift.
    \item[Mean time to detection (MTD):] the mean number of samples between a drift point and its first matching detection, where a missed drift point counts as $T$.
    \item[False alarm:] a detection that is not matched to a drift point. On streams without drift every detection is a false alarm, and false alarms are reported per 10{,}000 samples.
\end{description}

\section{Research Contributions}
This thesis makes the following contributions:

\begin{enumerate}
    \item We propose a corrected classifier two-sample drift detector for streams, EADD, that combines a reservoir reference, a permutation test whose null comes from the stream's own history, a sequential permutation test (ASPT) and TreeSHAP attribution.
    \item Evidence that a reservoir reference finds slow feature drift that a sliding reference misses (recall 1.00 vs.\ 0.53 for D3-LightGBM), at no loss on abrupt drift.
    \item Evidence that the history null keeps false alarms low on autocorrelated streams (about 3.5 times fewer than D3-LightGBM), and a measured account of the test's validity.
    \item Evidence that ASPT keeps detection quality at a third of the permutations.
    \item An empirical check, on data streams with known drifting features, that TreeSHAP attribution in a classifier-based drift detector names the drifting features (measured with EADD: the top-ranked feature drifted in 98\% of detections).
    \item A registered, reproducible evaluation, with the code, the raw runs and the analysis plans public.
\end{enumerate}

\section{Organisation of the Thesis}
Chapter~\ref{chap:literature} reviews drift types, unsupervised drift detectors grouped by approach, permutation and sequential tests, and feature attribution for drift. Chapter~\ref{chap:methodology} describes EADD, the baseline detectors, the data streams, the metrics and the registered experimental design. Chapter~\ref{chap:results} reports the results by hypothesis. Chapter~\ref{chap:discussion} interprets them and states the limitations of the study, and Chapter~\ref{chap:conclusion} concludes and recommends further work.
